% Options for packages loaded elsewhere
\PassOptionsToPackage{unicode}{hyperref}
\PassOptionsToPackage{hyphens}{url}
\PassOptionsToPackage{dvipsnames,svgnames,x11names}{xcolor}
\documentclass[
  11pt,
  letterpaper,
]{article}
\usepackage{xcolor}
\usepackage[margin=1in]{geometry}
\usepackage{amsmath,amssymb}
\setcounter{secnumdepth}{-\maxdimen} % remove section numbering
\usepackage{iftex}
\ifPDFTeX
  \usepackage[T1]{fontenc}
  \usepackage[utf8]{inputenc}
  \usepackage{textcomp} % provide euro and other symbols
\else % if luatex or xetex
  \usepackage{unicode-math} % this also loads fontspec
  \defaultfontfeatures{Scale=MatchLowercase}
  \defaultfontfeatures[\rmfamily]{Ligatures=TeX,Scale=1}
\fi
\usepackage{lmodern}
\ifPDFTeX\else
  % xetex/luatex font selection
\fi
% Use upquote if available, for straight quotes in verbatim environments
\IfFileExists{upquote.sty}{\usepackage{upquote}}{}
\IfFileExists{microtype.sty}{% use microtype if available
  \usepackage[]{microtype}
  \UseMicrotypeSet[protrusion]{basicmath} % disable protrusion for tt fonts
}{}
\usepackage{setspace}
\makeatletter
\@ifundefined{KOMAClassName}{% if non-KOMA class
  \IfFileExists{parskip.sty}{%
    \usepackage{parskip}
  }{% else
    \setlength{\parindent}{0pt}
    \setlength{\parskip}{6pt plus 2pt minus 1pt}}
}{% if KOMA class
  \KOMAoptions{parskip=half}}
\makeatother
\usepackage{color}
\usepackage{fancyvrb}
\newcommand{\VerbBar}{|}
\newcommand{\VERB}{\Verb[commandchars=\\\{\}]}
\DefineVerbatimEnvironment{Highlighting}{Verbatim}{commandchars=\\\{\}}
% Add ',fontsize=\small' for more characters per line
\newenvironment{Shaded}{}{}
\newcommand{\AlertTok}[1]{\textcolor[rgb]{1.00,0.00,0.00}{\textbf{#1}}}
\newcommand{\AnnotationTok}[1]{\textcolor[rgb]{0.38,0.63,0.69}{\textbf{\textit{#1}}}}
\newcommand{\AttributeTok}[1]{\textcolor[rgb]{0.49,0.56,0.16}{#1}}
\newcommand{\BaseNTok}[1]{\textcolor[rgb]{0.25,0.63,0.44}{#1}}
\newcommand{\BuiltInTok}[1]{\textcolor[rgb]{0.00,0.50,0.00}{#1}}
\newcommand{\CharTok}[1]{\textcolor[rgb]{0.25,0.44,0.63}{#1}}
\newcommand{\CommentTok}[1]{\textcolor[rgb]{0.38,0.63,0.69}{\textit{#1}}}
\newcommand{\CommentVarTok}[1]{\textcolor[rgb]{0.38,0.63,0.69}{\textbf{\textit{#1}}}}
\newcommand{\ConstantTok}[1]{\textcolor[rgb]{0.53,0.00,0.00}{#1}}
\newcommand{\ControlFlowTok}[1]{\textcolor[rgb]{0.00,0.44,0.13}{\textbf{#1}}}
\newcommand{\DataTypeTok}[1]{\textcolor[rgb]{0.56,0.13,0.00}{#1}}
\newcommand{\DecValTok}[1]{\textcolor[rgb]{0.25,0.63,0.44}{#1}}
\newcommand{\DocumentationTok}[1]{\textcolor[rgb]{0.73,0.13,0.13}{\textit{#1}}}
\newcommand{\ErrorTok}[1]{\textcolor[rgb]{1.00,0.00,0.00}{\textbf{#1}}}
\newcommand{\ExtensionTok}[1]{#1}
\newcommand{\FloatTok}[1]{\textcolor[rgb]{0.25,0.63,0.44}{#1}}
\newcommand{\FunctionTok}[1]{\textcolor[rgb]{0.02,0.16,0.49}{#1}}
\newcommand{\ImportTok}[1]{\textcolor[rgb]{0.00,0.50,0.00}{\textbf{#1}}}
\newcommand{\InformationTok}[1]{\textcolor[rgb]{0.38,0.63,0.69}{\textbf{\textit{#1}}}}
\newcommand{\KeywordTok}[1]{\textcolor[rgb]{0.00,0.44,0.13}{\textbf{#1}}}
\newcommand{\NormalTok}[1]{#1}
\newcommand{\OperatorTok}[1]{\textcolor[rgb]{0.40,0.40,0.40}{#1}}
\newcommand{\OtherTok}[1]{\textcolor[rgb]{0.00,0.44,0.13}{#1}}
\newcommand{\PreprocessorTok}[1]{\textcolor[rgb]{0.74,0.48,0.00}{#1}}
\newcommand{\RegionMarkerTok}[1]{#1}
\newcommand{\SpecialCharTok}[1]{\textcolor[rgb]{0.25,0.44,0.63}{#1}}
\newcommand{\SpecialStringTok}[1]{\textcolor[rgb]{0.73,0.40,0.53}{#1}}
\newcommand{\StringTok}[1]{\textcolor[rgb]{0.25,0.44,0.63}{#1}}
\newcommand{\VariableTok}[1]{\textcolor[rgb]{0.10,0.09,0.49}{#1}}
\newcommand{\VerbatimStringTok}[1]{\textcolor[rgb]{0.25,0.44,0.63}{#1}}
\newcommand{\WarningTok}[1]{\textcolor[rgb]{0.38,0.63,0.69}{\textbf{\textit{#1}}}}
\usepackage{longtable,booktabs,array}
\newcounter{none} % for unnumbered tables
\usepackage{calc} % for calculating minipage widths
% Correct order of tables after \paragraph or \subparagraph
\usepackage{etoolbox}
\makeatletter
\patchcmd\longtable{\par}{\if@noskipsec\mbox{}\fi\par}{}{}
\makeatother
% Allow footnotes in longtable head/foot
\IfFileExists{footnotehyper.sty}{\usepackage{footnotehyper}}{\usepackage{footnote}}
\makesavenoteenv{longtable}
\usepackage{graphicx}
\makeatletter
\newsavebox\pandoc@box
\newcommand*\pandocbounded[1]{% scales image to fit in text height/width
  \sbox\pandoc@box{#1}%
  \Gscale@div\@tempa{\textheight}{\dimexpr\ht\pandoc@box+\dp\pandoc@box\relax}%
  \Gscale@div\@tempb{\linewidth}{\wd\pandoc@box}%
  \ifdim\@tempb\p@<\@tempa\p@\let\@tempa\@tempb\fi% select the smaller of both
  \ifdim\@tempa\p@<\p@\scalebox{\@tempa}{\usebox\pandoc@box}%
  \else\usebox{\pandoc@box}%
  \fi%
}
% Set default figure placement to htbp
\def\fps@figure{htbp}
\makeatother
\setlength{\emergencystretch}{3em} % prevent overfull lines
\providecommand{\tightlist}{%
  \setlength{\itemsep}{0pt}\setlength{\parskip}{0pt}}
% header.tex — preamble additions for the lualatex build.
% Goal: an academic Latin Modern look, with a belt-and-suspenders guarantee that
% none of the 20 non-ASCII glyphs in paper.md can silently drop.

% (fontspec is already loaded by pandoc's template for lualatex; re-stating is harmless.)
% Academic Latin Modern body. Non-ASCII coverage is guaranteed by the \newunicodechar
% net below (the luaotfload font-fallback feature is unreliable on this MiKTeX build,
% so we do not use it — every non-ASCII glyph in paper.md is mapped explicitly).
\usepackage{fontspec}
\setmainfont{Latin Modern Roman}
\setsansfont{Latin Modern Sans}
\setmonofont{Latin Modern Mono}

\usepackage{microtype}
\usepackage{xurl}            % break long URLs (dashboard/github links) instead of overflowing
\setlength{\emergencystretch}{3em}

% Glyph safety net: map the math-ish Unicode used in prose to real LaTeX glyphs,
% so they render even if the font fallback misses (independent guarantee).
\usepackage{newunicodechar}
\newunicodechar{×}{\ensuremath{\times}}
\newunicodechar{→}{\ensuremath{\rightarrow}}
\newunicodechar{⇒}{\ensuremath{\Rightarrow}}
\newunicodechar{≈}{\ensuremath{\approx}}
\newunicodechar{≤}{\ensuremath{\leq}}
\newunicodechar{−}{\ensuremath{-}}
\newunicodechar{∈}{\ensuremath{\in}}
\newunicodechar{∎}{\ensuremath{\blacksquare}}
\newunicodechar{·}{\ensuremath{\cdot}}
\newunicodechar{…}{\ldots}
\newunicodechar{⁵}{\textsuperscript{5}}
\newunicodechar{α}{\ensuremath{\alpha}}
\newunicodechar{γ}{\ensuremath{\gamma}}
\newunicodechar{τ}{\ensuremath{\tau}}
\newunicodechar{ε}{\ensuremath{\varepsilon}}
\newunicodechar{Γ}{\ensuremath{\Gamma}}
% (§, em/en dash, middle interpunct, ï are covered natively by Latin Modern Roman.)
\usepackage{bookmark}
\IfFileExists{xurl.sty}{\usepackage{xurl}}{} % add URL line breaks if available
\urlstyle{same}
% fallback for those not using the hyperref driver hyperxmp:
\makeatletter
\@ifundefined{xmpquote}{\newcommand{\xmpquote}[1]{#1}}{}
\makeatother
\hypersetup{
  pdftitle={Overlapping{,} Not Approximating: Asynchronous-Style Speedup in Reasoning-RL Training and the Lossless Design of TLT},
  pdfauthor={Eric Hibbs, independent analysis (not the original ASPLOS paper); see provenance note},
  colorlinks=true,
  linkcolor={Maroon},
  filecolor={Maroon},
  citecolor={Blue},
  urlcolor={Blue},
  pdfcreator={LaTeX via pandoc}}

\title{Overlapping, Not Approximating: Asynchronous-Style Speedup in
Reasoning-RL Training and the Lossless Design of TLT}
\author{Eric Hibbs, independent analysis (not the original ASPLOS
paper); see provenance note}
\date{Built 2026-06-08 from paper.md (frozen content; mtime 2026-06-08
09:07)}

\begin{document}
\maketitle

% frontmatter.tex — inserted after \maketitle, before the TOC.
% Provenance framing + four-axis legend (metadata framing, not content). No bold, no em-dashes.
\begin{center}
\fbox{\parbox{0.93\linewidth}{\small
\textit{Independent analysis, not the original ASPLOS paper.}\;
This document is an independent analysis and synthesis built on Hu, Yang \emph{et al.},
\emph{Taming the Long-Tail: Efficient Reasoning RL Training with Adaptive Drafter}
(ASPLOS~'26, arXiv:2511.16665v2, CC-BY-4.0). Analytical models are adapted from the interactive
companion dashboard by Eric Hibbs, and mechanism claims are grounded in the released
code (\texttt{mit-han-lab/fastrl}, Apache-2.0). All original experimental results belong to
the paper's authors.

\medskip
\textit{Provenance legend (four axes).}\;
\texttt{[Paper:\ §/Fig/Tbl]} = measured in the source paper;\;
\texttt{[Dashboard:\ Model Math §N]} = modeled, assumption-based (Hibbs), not a paper result;\;
\texttt{[Code:\ path]} = confirmed in the released fastrl source;\;
\texttt{[Author synthesis]} = this document's own reasoning.
}}
\end{center}
\vspace{0.6em}

{
\setcounter{tocdepth}{3}
\tableofcontents
}
\listoffigures
\listoftables
\setstretch{1.05}
\subsection{Abstract}\label{abstract}

Reinforcement-learning training of reasoning LLMs is throttled by a
\emph{long-tail} rollout problem: a handful of very long generations
dominate each step while most workers sit idle. TLT attacks this not by
approximating the algorithm but by overlapping otherwise-idle work in
time. Its two pillars, an \emph{Adaptive Drafter} trained
opportunistically on idle GPUs during the tail, and an \emph{Adaptive
Rollout Engine} that turns on speculative decoding (SD) exactly when the
batch shrinks into SD's sweet spot, convert wasted tail bubbles into
useful compute while keeping training token-level lossless and strictly
on-policy. The paper reports ``over 1.7×'' end-to-end speedup over
state-of-the-art systems {[}Paper Abstract{]}, with the figure-level
range given as 1.7--2.1× over VeRL {[}Paper Fig 11{]}, a geomean of
1.86× on H100 and 1.70× on A100, and a peak of 2.12× on Qwen-32B/H100
{[}Paper Fig 11, Tbl 3{]}. Rollout alone is ≈85\% of step time {[}Paper
§1, Fig 1a{]}. This paper frames TLT around its
\emph{asynchrony/concurrency} mechanisms, non-blocking drafter training,
preemptible opportunistic scheduling, selective asynchronous
checkpointing, and tail-phase adaptive SD, pairs each with its measured
speedup \emph{and its implementation in the released code}, develops a
mathematical model of where the gains come from and why they saturate
(using the dashboard's order-statistic, Markov-acceptance,
optimal-stopping, and moving-target-convergence derivations, each with
its assumptions flagged), and contrasts TLT's \emph{lossless} synchrony
with the off-policy drift risked by fully asynchronous RL systems. The
thesis: you can capture async-style utilization gains without paying the
correctness tax of async RL.

\begin{center}\rule{0.5\linewidth}{0.5pt}\end{center}

\begin{quote}
Provenance conventions (four axes). Every quantitative or mechanistic
claim is tagged: {[}Paper §/Fig/Tbl{]} = stated in the TLT paper
(\emph{measured}); {[}Dashboard §N{]} = derived/assumed by the dashboard
author (Hibbs) in ``Model Math'' (\emph{modeled, assumption-based, NOT a
paper result}); {[}Code: path{]} = confirmed against the released repo
(\emph{implemented}; \texttt{mit-han-lab/fastrl} @ commit
\texttt{bce3df7}); {[}Author synthesis{]} = my own derivation or framing
(\emph{synthesized}). Direct quotes are verbatim, \textless{} 15 words,
at most one per source per claim; everything else is paraphrase.
``Lossless'' throughout means token-level distributional equivalence of
rollouts, speculative-decoding verification preserves the target model's
exact sampling distribution; it is not a claim about convergence or
final accuracy. Throughput numbers are per-step token throughput at
matched reward curves, not wall-clock-to-convergence.
\end{quote}

\subsection{1. Introduction}\label{introduction}

\subsubsection{1.1 The long-tail
bottleneck}\label{the-long-tail-bottleneck}

A GRPO-style RL step has three stages, rollout (generate candidate
responses), inference (score them and compute a KL penalty), and
training (update the policy). The first stage dominates: the paper
measures the rollout phase at approximately 85\% of total step time
{[}Paper §1 item (1), Fig 1a{]}. Worse, rollout latency is governed by
the \emph{slowest} response in the batch, and response lengths are
heavy-tailed: most are short, a few hit the configured maximum. The
production evidence is stark. A ByteDance trace {[}Paper Fig 2{]} of
Qwen2.5-32B on 128 H20-96GB GPUs ran 11 days for only 385 steps,
averaging ≈40 min/step, with evaluation every 5 steps costing ≈20 min
each, under a customized max length of 20,480 tokens {[}Paper §1 item
(3), Fig 2{]}. Because the batch finishes only when its longest member
finishes, the gap between the maximum and the 75th percentile is a
direct measure of wasted GPU-time, an ``under-utilized zone'' visible
across the entire run {[}Paper Fig 2{]}. (A schematic of this
distribution and the under-utilized zone is in
\hyperref[figures]{Figures}, \texttt{fig1\_longtail\_SCHEMATIC},
illustrative only, since the raw per-request distribution is not
published as numbers.)

\subsubsection{\texorpdfstring{1.2 Why asynchronous \emph{overlap}, not
approximation}{1.2 Why asynchronous overlap, not approximation}}\label{why-asynchronous-overlap-not-approximation}

There are two ways to make a tail-bound workload faster. One is to
approximate: drop the stragglers, accept stale gradients, or mix policy
versions, the route taken by fully asynchronous RL. These trade
correctness for utilization. The other is to overlap: keep the algorithm
exact, but fill the idle time the tail creates with useful, independent
work. TLT takes the second route. Its design rests on two insights
stated by the paper: rollout bubbles can be exploited because resources
freed as short sequences finish can be repurposed at no extra cost, and
drafter training can be decoupled from rollout completion so that its
computation overlaps ongoing rollouts rather than following them
{[}Paper §1, §4.2{]}. Speculative decoding is the ideal payload because
it is ``Mathematically Lossless'' {[}Paper §2.2{]}, the accepted token
distribution is identical to the target's, and it is most effective
precisely in the tail, where effective batch sizes are small and
decoding is memory-bound. The result is \emph{async-style utilization}
(idle GPUs doing work) with \emph{synchronous-grade correctness} (the RL
update is unchanged).

\subsubsection{1.3 Contributions of this
paper}\label{contributions-of-this-paper}

This is a survey-and-analysis paper, not a new system. Its contributions
are: (i) a reframing of TLT around the concurrency mechanisms that
produce its speedups, each paired with its measured number and its
location in the released code (§4, grounded by a full mechanism→code
audit); (ii) a consolidated mathematical model of the speedup and its
saturation, built from the dashboard's derivations with every assumption
flagged and connected to the paper's tables (§3, §5); (iii) a synthesis
of TLT's evaluation (§6); (iv) an analysis of why TLT stays losslessly
synchronous where async-RL systems risk off-policy drift (§7); (v) an
Implementation \& Reproducibility account and a Threats-to-Validity /
Assumptions ledger (§8--§9); and (vi) a four-way-attributed treatment
(paper / dashboard / code / synthesis) of every claim.

\begin{center}\rule{0.5\linewidth}{0.5pt}\end{center}

\subsection{2. Background and
Motivation}\label{background-and-motivation}

\subsubsection{2.1 Reasoning RL and GRPO's three
stages}\label{reasoning-rl-and-grpos-three-stages}

Reasoning LLMs improve through \emph{test-time scaling}: longer chains
of thought and self-reflection raise accuracy, and these behaviors are
largely acquired during RL {[}Paper §2.1{]}. The paper uses Group
Relative Policy Optimization (GRPO), which replaces PPO's learned critic
with a group-relative baseline, cutting overhead {[}Paper §2.1{]}. A
single GRPO step {[}Paper Fig 4{]} is:

\begin{enumerate}
\def\labelenumi{\arabic{enumi}.}
\tightlist
\item
  Rollout, the target model generates multiple responses per prompt
  (autoregressive decoding; the long-tail stage). The released runs use
  a group size of n = 8 responses/prompt {[}Code:
  examples/grpo\_7B.sh{]}.
\item
  Inference, responses are prefilled by both the target and a frozen
  reference model to compute a KL-divergence penalty; a rule-based
  reward (e.g., answer correctness) is computed. The code applies KL as
  a loss term (\texttt{low\_var\_kl}, coefficient 0.001), not folded
  into the reward {[}Code: examples/grpo\_7B.sh{]}.
\item
  Training, a loss combining reward and KL updates the target weights.
\end{enumerate}

The asymmetry matters: rollout is sequential token-by-token generation
(latency-bound, tail-dominated), while inference and training are
batched prefills/backward passes. TLT inserts its opportunistic work
into stage 1's tail and reuses stage 2's hidden states.

\subsubsection{2.2 Speculative decoding: lossless, and a memory→compute
shift}\label{speculative-decoding-lossless-and-a-memorycompute-shift}

In speculative decoding a small draft model proposes a short run of
tokens, and the large target verifies them all in one parallel forward
pass, accepting the longest correct prefix plus one corrected token
{[}Paper §2.2, Fig 5{]}. Two properties make SD ideal for RL. First, it
is token-level lossless: the verification step preserves the target's
exact next-token sampling distribution, so the accepted sequence is
distributed identically to ordinary autoregressive sampling from the
target (this is a \emph{distributional} guarantee on rollouts, not a
statement about convergence). Second, it changes the \emph{arithmetic
intensity}: autoregressive decoding at small batch is memory-bound (one
token's weights streamed per step); verifying γ speculative tokens at
once is compute-bound, better saturating the GPU. The paper's roofline
analysis shows SD reaching peak TFLOPS at far smaller batch sizes than
vanilla decoding {[}Paper Fig 5c{]}, i.e., SD's benefit is largest
exactly where the tail lives (few in-flight requests). The paper
summarizes the mechanism as shifting from ``memory-bound towards
compute-bound'' {[}Paper §2.2{]}.

Three challenges block naïve SD in RL {[}Paper §2.2/§2
``Challenges''{]}: C1, the evolving target model makes a fixed drafter
go stale; C2, drafter training cost; C3, fluctuating batch sizes, since
SD that helps at batch 1 can hurt or OOM at large batch. TLT's two
components answer C1/C2 (Adaptive Drafter) and C3 (Adaptive Rollout
Engine).

\subsubsection{2.3 What ``asynchronous'' means
here}\label{what-asynchronous-means-here}

This is the crux of the framing, and the word is overloaded. In async-RL
systems (e.g., AReaL), ``asynchronous'' means \emph{algorithmic}
asynchrony: the policy is updated from rollouts generated by
\emph{earlier} policy versions, so training and generation run on
different model versions concurrently, fast, but off-policy. In TLT,
``asynchronous'' means \emph{systems-level} overlap with \emph{no}
algorithmic change: the policy update still consumes only on-policy
rollouts of the current version; what runs concurrently is the
\emph{drafter's} training and the \emph{checkpoint's} save, neither of
which touches the policy gradient. TLT's first design principle is an
explicit ``Lossless Guarantee'' that forbids asynchronous staleness in
the RL update itself {[}Paper §3{]}. So the entire speedup must come
from overlapping \emph{auxiliary} work, never from desynchronizing the
\emph{core} work, the distinction this paper builds on.

\begin{center}\rule{0.5\linewidth}{0.5pt}\end{center}

\subsection{3. The Long-Tail, Formally}\label{the-long-tail-formally}

The paper's evidence for the tail is empirical: the response-length
histogram with a max length of 30K for Qwen-7B and DeepSeek-R1-7B
{[}Paper Fig 1a{]}, and the ByteDance percentile trace {[}Paper Fig
2{]}. The dashboard supplies a \emph{motivating} closed-form analysis on
top of this. It is useful intuition but rests on an assumption the paper
does not make, so it is flagged throughout.

\subsubsection{3.1 Batch completion as an order
statistic}\label{batch-completion-as-an-order-statistic}

Let a batch hold \(B\) responses with token-lengths \(L_1,\dots,L_B\).
If response \(i\) decodes at throughput \(\text{throughput}_i\), the
batch finishes only when its slowest member does:
\[T_{\text{batch}} = \max_{i=1}^{B}\frac{L_i}{\text{throughput}_i}.\qquad\text{[Dashboard §1]}\]
The bottleneck is therefore the maximum, a classical order statistic,
the formal version of ``a few very long responses dominate execution
time'' {[}Paper Abstract{]}.

\subsubsection{3.2 Expected maximum under a Pareto tail (with the
assumption
flagged)}\label{expected-maximum-under-a-pareto-tail-with-the-assumption-flagged}

The dashboard models the tail as Pareto,
\(P(L>x) = (x_{\min}/x)^{\alpha}\) with shape \(\alpha>1\), i.e.~CDF
\(F(x)=1-(x_{\min}/x)^{\alpha}\) {[}Dashboard §1{]}. This is an
{[}Assumption{]}: the dashboard explicitly notes the paper's Fig 1a is
empirical (log-normal-shaped) and that Pareto is chosen only for a
tractable closed form {[}Dashboard §1 assumption tag{]}. With that
caveat, the dashboard states the expected maximum as
\[\mathbb{E}[M_B] = x_{\min}\cdot B^{1/\alpha}\cdot\frac{\alpha}{\alpha-1}.\qquad\text{[Dashboard §1]}\]

Step-by-step derivation {[}Author synthesis{]}. Let \(M_B=\max_i L_i\)
with the \(L_i\) i.i.d. Pareto. Then \(P(M_B\le x)=F(x)^B\). The exact
expectation is
\[\mathbb{E}[M_B] = x_{\min}\,\frac{\Gamma(B+1)\,\Gamma\!\left(1-\tfrac{1}{\alpha}\right)}{\Gamma\!\left(B+1-\tfrac{1}{\alpha}\right)},\qquad \alpha>1.\]
Using the ratio asymptotic
\(\Gamma(B+1)/\Gamma(B+1-\tfrac{1}{\alpha})\sim B^{1/\alpha}\) for large
\(B\) (Stirling), this becomes
\(\mathbb{E}[M_B]\approx x_{\min}B^{1/\alpha}\Gamma(1-\tfrac1\alpha)\).
The dashboard's form replaces the tail constant
\(\Gamma(1-\tfrac1\alpha)\) with \(\tfrac{\alpha}{\alpha-1}\), the
Pareto mean \(\mathbb{E}[L]=x_{\min}\tfrac{\alpha}{\alpha-1}\). So the
dashboard's closed form is best read as \emph{``scale the mean by the
order-statistic growth factor \(B^{1/\alpha}\).''}

What it means physically. The straggler grows as \(B^{1/\alpha}\): the
heavier the tail (smaller \(\alpha\)), the faster the worst case blows
up as you widen the batch, and the deeper the rollout bubble, and, read
in reverse, the more freed capacity becomes available for opportunistic
work.

Numerical recomputation and a flagged error {[}Author synthesis{]}. I
recomputed the α=2, B=128 case (exact Γ formula + 2×10⁵-sample
Monte-Carlo): \(\mathbb{E}[M_B]\approx\mathbf{20.07\,x_{\min}}\) (exact)
/ 20.14 (MC); the dashboard's own closed form gives 22.63 (it mildly
over-states, since \(\tfrac{\alpha}{\alpha-1}=2\) exceeds
\(\Gamma(\tfrac12)=\sqrt\pi\approx1.77\)). The Pareto median is
\(x_{\min}\sqrt2\), so the true expected-max/median ratio is ≈14×
(exact/MC), or 16× under the dashboard's formula. The dashboard's worked
example states ``≈5.7× vs median'' {[}Dashboard §1{]}, this is
numerically wrong: \(5.7\approx\sqrt{128}/2=B^{1/\alpha}/\alpha\), a
mis-normalization (it drops the mean factor and mis-divides). The
qualitative point (the straggler is many-fold the median) holds, but the
\emph{specific} multiplier is the dashboard's and should not be cited as
a derived result; the correct figure is ≈14×.

\subsubsection{3.3 From the math back to the paper's
measurement}\label{from-the-math-back-to-the-papers-measurement}

The order-statistic view predicts that under-utilization is structural,
not incidental, which is exactly what the paper observes empirically and
persistently across 385 steps {[}Paper Fig 2{]}, and what motivates
harvesting the bubble rather than merely shortening it.

\begin{center}\rule{0.5\linewidth}{0.5pt}\end{center}

\subsection{4. Mechanisms of Async
Speedup}\label{mechanisms-of-async-speedup}

Each subsection names a concurrency mechanism, pairs it with its
measured effect {[}Paper{]}, and grounds it in the released code
{[}Code{]} (full audit: \texttt{phase1\_code\_grounding.md}; all 11
mechanisms located and confirmed). Collectively they convert the tail
bubble (§3) into the end-to-end gains of §6.

\subsubsection{4.1 Non-blocking Spot Trainer + DataBuffer (overlapping
drafter training with
rollout)}\label{non-blocking-spot-trainer-databuffer-overlapping-drafter-training-with-rollout}

The core asynchrony is decoupling drafter training from rollout
completion. Vanilla RL runs rollout → inference → training in sequence;
a naïve SpecRL bolts drafter training on as another sequential step
{[}Paper Fig 8a,b{]}. The Spot Trainer, class
\texttt{EagleBackgroundTrainer} {[}Code:
verl/workers/drafter/eagle\_background\_trainer.py:35{]}, instead trains
the drafter on \emph{partial} responses as soon as enough short
sequences finish, executing ``in a non-blocking manner'' during the tail
{[}Paper §4.2, Fig 8c{]}.

Training only on early finishers would bias the drafter away from long
sequences, so TLT adds the DataBuffer {[}Code:
verl/utils/data\_buffer.py{]}: it caches hidden states and tokens from
the inference stage and uses one-step-offset sampling, confirmed in code
as collecting data the step \emph{before} training
(\texttt{should\_collect\_data\_this\_step} fires at
\texttt{(step+1)\ \%\ interval\ ==\ 0}) and drawing the training batch
from the last \texttt{buffer\_steps\ =\ 2} RL steps {[}Code:
worker\_manager.py:490-493; eagle\_background\_trainer.py:276{]}.
\emph{Provenance nuance:} the paper describes prioritizing the previous
step's long sequences; the released code samples the last two steps via
\texttt{random.sample} over a response-centered 512-token window, the
cross-step offset is confirmed, but an explicit \emph{length-ranked}
prioritization is not present as written {[}Code; see §9{]}. The paper
reports that partial current-step data combined with buffered sequences
attains performance comparable to full-set training {[}Paper §4.2{]}.
The hidden states are \emph{reused} from the inference phase, so drafter
training adds no extra prefill; the target is ``over 20×'' larger than
the draft model {[}Paper §4.1{]}.

\emph{Measured effect:} this mechanism makes the drafter a ``free
byproduct'' {[}Paper Abstract{]}, its training cost is hidden inside the
bubble.

\subsubsection{4.2 Worker Coordinator state machine (preemptible
opportunistic
scheduling)}\label{worker-coordinator-state-machine-preemptible-opportunistic-scheduling}

The bubble must be harvested without disturbing rollout. A centralized
Worker Coordinator (\texttt{RolloutDrafterManager} +
\texttt{CentralCoordinator}, rank-0, ZeroMQ; REP/PUB on the coordinator,
REQ/SUB on workers) tracks each worker, one rollout instance, e.g.~a
TP=8 worker spans 8 GPUs {[}Paper §4.2; Code:
worker\_manager.py:91-107,441{]}. The paper summarizes the lifecycle as
BUSY → IDLE → TRAINING {[}Paper §4.2{]}; the released code implements a
five-state machine,
\texttt{IDLE\ /\ GENERATING\ /\ RELEASED\ /\ TRAINING\ /\ COMPLETED},
where the paper's ``BUSY'' = \texttt{GENERATING} and ``IDLE (memory
released)'' = \texttt{RELEASED} {[}Code: worker\_manager.py:54-61{]}.
Once enough workers release, the coordinator promotes them to training
(default threshold \texttt{min\_workers\_for\_training\ =\ 1}) and, on
rollout completion, broadcasts
\texttt{STOP\_TRAINING}/\texttt{BATCH\_COMPLETE} for graceful shutdown
{[}Code: worker\_manager.py:245-283,204-230; fastrl\_trainer.yaml:25{]}.
This realizes TLT's ``Interference-Free'' principle {[}Paper §3{]},
drafter work runs as low-priority ``spot tasks'' {[}Paper §4.2{]}. The
shipped drafter-training cadence is
\texttt{training\_interval\_steps\ =\ 10}, matching the paper's ``every
\textasciitilde10 steps is sufficient'' {[}Paper §6.4; Code:
fastrl\_trainer.yaml:20{]}.

\emph{Measured effect:} stage-transition overhead is \textless1\% of the
RL step {[}Paper §6.4{]}.

\subsubsection{4.3 Selective asynchronous checkpointing (overlapping the
save with
training)}\label{selective-asynchronous-checkpointing-overlapping-the-save-with-training}

Preemptibility implies frequent checkpoints, which would themselves
stall training. TLT offloads the save to a background thread
(\texttt{dcp.async\_save}, PyTorch DCP) and dumps only the drafter's
trainable parameters, excluding the frozen set
\texttt{\{model.embed\_tokens.weight,\ lm\_head.weight\}} {[}Code:
eagle\_background\_trainer.py:83,105-148{]}. The paper's measured ladder
{[}Paper Fig 17a{]}:

\begin{itemize}
\tightlist
\item
  Vanilla synchronous checkpoint: 893 ms → Async (background thread):
  280 ms (3.2×) → Selective async (background + trainable-only): 97 ms
  (9.2× vs vanilla) {[}Paper Fig 17a, §6.4{]}.
\end{itemize}

\emph{Reading the 9.2× honestly:} this is the paper's end-to-end
measurement, not a parameter-count ratio. The 3.2× comes from
overlapping I/O; the further step to 9.2× comes from selectivity, and
because the object being checkpointed is the single-layer drafter (where
the frozen, tied embedding + LM head dominate the parameter count),
excluding them is a large reduction. The async + selective
\emph{mechanism} is code-confirmed; the exact magnitude is the paper's
{[}Paper Fig 17a{]}.

\subsubsection{4.4 Drafter architecture and sequence packing (cheap
payload, dense
compute)}\label{drafter-architecture-and-sequence-packing-cheap-payload-dense-compute}

The drafter is a single trainable transformer decoder layer
(\texttt{config.num\_hidden\_layers\ =\ 1}) that reuses the target's
frozen Embedding and LM Head (\texttt{requires\_grad\ =\ False})
{[}Paper §4.1, Fig 7; Code:
eagle-train/eagle\_trainer.py:190,255-258{]}. Only that one layer
updates, ≈ 1/layer\_num of target parameters {[}Paper §4.1{]}. The
paper's contrast: Qwen2.5-32B has 64 layers vs Qwen2.5-0.5B's 24; a 64×
parameter cut still leaves latency dominated by sequential layers, so
TLT's single-layer drafter, at roughly the parameter count of a
Qwen2.5-0.5B drafter, is 2.4× faster {[}Paper §4.1{]}. \emph{Code
refinements:} the drafter loss is \texttt{SmoothL1Loss} (Huber) on
hidden states + a soft-target cross-entropy on logits, weighted 0.5/0.5,
a sharper statement of the paper's ``L1 + CE'' {[}Code:
eagle\_background\_trainer.py:73,558-578{]}; and the ``frozen LM head''
holds on the default EAGLE path but not on the EAGLE-3 path, where the
LM head is trainable for the draft vocabulary {[}Code:
eagle3\_trainer.py:475; see §9{]}.

Within the preemptible slots, sequence packing concatenates
variable-length sequences (attention-masked, no padding) to keep the GPU
dense {[}Paper §4.2; Code: eagle\_background\_trainer.py:410-417;
verl/models/mcore/util.py{]}. \emph{Measured effect:} training
throughput rises 13.3 → 29.6 samples/s (2.2×) over padded batching
{[}Paper Fig 17b{]}. (See \hyperref[figures]{Figures}:
\texttt{fig17\_spot\_trainer}.)

\subsubsection{4.5 Adaptive SD + Bucketed CUDAGraph + BEG-MAB
(tail-phase, OOM-safe,
self-tuning)}\label{adaptive-sd-bucketed-cudagraph-beg-mab-tail-phase-oom-safe-self-tuning}

SD must be turned on only where it helps. The Adaptive SD Manager
activates SD only when the remaining request count drops below a
threshold (default 32) {[}Paper §5.1; Code: examples/grpo\_7B.sh:28
\texttt{speculative.bs\_threshold=32}{]}. The code adds a hysteresis the
paper omits: SD enables only after the batch stays \texttt{≤\ 32} for 10
consecutive checks (\texttt{adaptive\_spec\_warmup\_checks}), and the
counter resets above threshold {[}Code: eagle\_worker.py:361,364,196{]}.
Tree-based drafting explores the topK branches for \texttt{Draft\_Depth}
steps and verifies the \texttt{Tokens\_to\_Verify} highest-confidence
tokens in parallel {[}Paper §5.1, Fig 9; Code:
eagle\_worker.py:748-821{]}.

Bucketed CUDAGraph Capture groups batch sizes into buckets,
disaggregates target vs draft graph capture, and merges identical
configurations {[}Paper §5.1; Code: server\_args.py:789-796{]}.
\emph{Measured effect} {[}Paper Tbl 5, Llama-3-8B, TP=4{]}: Single
strategy 7.81 GB; naïve multi-strategy 30.39 GB; Bucketed 10.69 GB, a
2.8× reduction vs naïve at near single-strategy cost.

Strategy selection is automated by the BEG-MAB tuner
(\texttt{BucketedEpsilonGreedyMAB}): each arm is
\texttt{(Draft\_Depth,\ topK,\ Tokens\_to\_Verify)}, bucketed by
\texttt{Tokens\_to\_Verify}, exploiting the median reward over a sliding
window {[}Paper §5.2, Alg 1; Code: eagle\_mab.py:197-292{]}. The exact
reward is code-confirmed:
\texttt{a\ =\ sum(accept\_len)/batch\_size\ +\ 1},
\texttt{reward\ =\ a\ ·\ batch\_size\ /\ elapsed\_time} {[}Code:
eagle\_worker.py:1234,1240{]}, matching Algorithm 1. Code-level defaults
the paper omits: ε = 0.1, window = 1000, bucket thresholds
{[}1,2,5,21{]}, and an OOM guard
\texttt{draft\_tokens\ ×\ batch\_size\ ≤\ 2048} {[}Code:
eagle\_mab.py:210,220,442{]}. For early steps, a model-free n-gram
drafter (Lookahead-style) is the fallback, with defaults
\texttt{max\_match\_window\ =\ 12}, \texttt{branch\_length\ =\ 18}
{[}Paper §5.3; Code: ngram\_worker.py; server\_args.py:342-347{]}.

\emph{Measured effect (case study)} {[}Paper Fig 14{]}: 128 requests on
Qwen-32B/H100, SD switches on below 32 remaining requests and delivers
an overall 2.44× over the no-SD baseline.

\begin{center}\rule{0.5\linewidth}{0.5pt}\end{center}

\subsection{5. A Mathematical Model of the
Speedup}\label{a-mathematical-model-of-the-speedup}

Every closed form here is the dashboard author's {[}Dashboard §2--§4{]},
\emph{modeled, assumption-based}; the paper's numbers are the empirical
anchors, and \texttt{phase3\_accuracy.md} records the numerical
re-derivations. These models are qualitative, not predictive fits to the
paper's tables.

\subsubsection{5.1 Speculative decoding as a truncated geometric
process}\label{speculative-decoding-as-a-truncated-geometric-process}

Assume the drafter proposes \(\gamma\) tokens and each is independently
accepted with probability \(\alpha\). This uniform per-token accept rate
is an {[}Assumption{]}, real rates decline with position (paper Fig 16)
{[}Dashboard §2{]}. Then
\[\mathbb{E}[n] = \frac{1-\alpha^{\gamma+1}}{1-\alpha}.\qquad\text{[Dashboard §2]}\]

Derivation {[}Author synthesis{]}. With \(A\) the accepted-prefix
length, \(P(A\ge k)=\alpha^{k}\); the per-position law is
\(P(A=k)=\alpha^k(1-\alpha)\) for \(k<\gamma\) and
\(P(A=\gamma)=\alpha^\gamma\). Tokens committed \(=A+1\), so
\(\mathbb{E}[n]=\sum_{k=0}^{\gamma}\alpha^k\); the geometric sum gives
\(\frac{1-\alpha^{\gamma+1}}{1-\alpha}\). ∎ Verified:
\(\alpha\to1\Rightarrow\gamma+1\) (computed 9.000 at γ=8);
\(\gamma\to\infty\Rightarrow\frac1{1-\alpha}\) (cap 5.0 at α=0.8). This
tapering matches the paper's accept-length saturation {[}Paper Fig
13a{]}, max 9.78 at depth 16/verify 128. Violation: Fig 16 shows α
decaying with depth, so the uniform-α premise does not hold in practice,
the model is \emph{qualitative}.

\subsubsection{5.2 The speedup, and why there is an optimal
depth}\label{the-speedup-and-why-there-is-an-optimal-depth}

With \(c=T_{\text{verify}}/T_{\text{draft}}\):
\[\text{speedup}=\frac{\mathbb{E}[n]}{1+\gamma/c}.\qquad\text{[Dashboard §2]}\]
Reading {[}Author synthesis{]}. Numerator saturates in \(\gamma\)
(§5.1); denominator grows linearly ⇒ an inverted-U with interior
optimum. Verified: \(\arg\max_\gamma\) speedup = 3/5/7 for \(c\) =
4/8/16, and the optimum grows with \(c\), qualitatively matching the
paper's grid peak of 3.66× at depth 12/verify 64 {[}Paper Fig 13b{]}.
The companion topK sweep is nearly flat, accept length 8.29--8.67,
speedup 3.47×--3.65× across topK 4--16 {[}Paper Tbl 1{]}, so TLT fixes
topK and tunes the rest. \emph{Note on conditions:} the SD grids (Fig
13/Tbl 1) are measured at temperature 0 for stable accept-length
readings; the operative accept rates during training (temp 0.9) are
those in Tbl 6/Fig 16. (See \hyperref[figures]{Figures}:
\texttt{fig13\_sd\_grid\_heatmaps},
\texttt{dashboard\_markov\_surface}.)

\subsubsection{5.3 The adaptive threshold as optimal stopping (why SD
collapses at large
batch)}\label{the-adaptive-threshold-as-optimal-stopping-why-sd-collapses-at-large-batch}

The dashboard models per-request SD benefit at \(k\) active requests as
exponential decay and asks when aggregate benefit clears the switch cost
\(O_{\text{SD}}\):
\[b(k)=b_0\,e^{-k/\tau},\qquad\text{activate when } k\cdot b(k)\ge O_{\text{SD}}.\qquad\text{[Dashboard §3]}\]
\(\tau=16\) is an {[}Assumption{]} fit to Table 4, not in the paper
{[}Dashboard §3{]}. Reading {[}Author synthesis{]}.
\(g(k)=b_0 k e^{-k/\tau}\) peaks at \(k=\tau\) then decays, so beyond
some \(k\) aggregate benefit no longer clears the overhead, an
activation threshold. \emph{Implementation reality:} the deployed rule
is not this smooth model but a hard \texttt{batch\_size\ ≤\ 32} step
with a 10-check warmup {[}Code: eagle\_worker.py:361,196{]}; the
\(b(k)\) model only explains the \emph{direction} of the collapse.
Empirical anchor, and a dashboard error. Paper Tbl 4 (Qwen-32B, TP=4,
depth 10, topK 8) shows, at verify-64, BS=1 → 3.62× falling to BS=32 →
1.70× (2.48× at verify-16). The dashboard's §3 text instead claims SD
``drops to 1.08× at BS=32'' while citing Table 4 {[}Dashboard §3{]}, but
no 1.08× appears in Table 4, and its companion ``3.66× at BS=1'' is the
Fig 13 grid peak, not a Tbl 4 entry {[}Paper Tbl 4; Fig 13b{]}. The
error is the dashboard's; only the paper's values are used here. (See
\hyperref[figures]{Figures}: \texttt{tbl4\_batch\_erosion},
\texttt{dashboard\_threshold\_curve}.)

\subsubsection{5.4 The drafter as a moving target (convergence under
policy
drift)}\label{the-drafter-as-a-moving-target-convergence-under-policy-drift}

The drafter minimizes
\(\mathcal{L}(\theta_t)=\mathbb{E}_{x\sim\pi_t}[-\log p_\theta(x)]\)
{[}Dashboard §4{]}. Using \(\alpha\approx e^{-D_{\text{KL}}}\) locally
and a per-step KL recursion (with \(\eta_d\) the drafter step gain,
\(I_t\) Fisher information, \(\delta_\pi\) policy drift):
\[D_{\text{KL}}(\pi_{t+1}\Vert p_{\theta_{t+1}})\le D_{\text{KL}}(\pi_t\Vert p_{\theta_t})-\eta_d I_t+\delta_\pi+O(\eta_d^2),\]
the accept rate is non-decreasing in expectation iff
\(\eta_d I_t>\delta_\pi\) {[}Dashboard §4{]}. Both rest on an
{[}Assumption{]} (Fisher ≈ constant; local KL--accept relation). What it
means {[}Author synthesis{]}. The drafter keeps up iff its per-step
learning outpaces policy drift, the skeleton of three paper
observations: the dip-then-recover after each target update {[}Paper Fig
15{]}; co-training lifting RL accept length 4.59→6.53 and downstream
3.76→5.15 {[}Paper Tbl 6{]}; and sustained accept rate at distant token
indices {[}Paper Fig 16{]}. Arithmetically \((6.53-4.59)/4.59=42.3\%\)
{[}Author synthesis{]}, matching the dashboard's ``+42\%''; the
dashboard reads Fig 16 as ≈32\% (adaptive) vs ≈6\% (vanilla) at index 8,
the dashboard's reading, not a printed paper value {[}Dashboard §4{]}.
(See \hyperref[figures]{Figures}: \texttt{tbl6\_adaptive\_drafter},
\texttt{fig16\_accept\_decay\_SCHEMATIC}.)

\begin{center}\rule{0.5\linewidth}{0.5pt}\end{center}

\subsection{6. Evaluation Synthesis}\label{evaluation-synthesis}

\subsubsection{6.1 Setup (so the numbers carry their
conditions)}\label{setup-so-the-numbers-carry-their-conditions}

TLT is built on VeRL + Ray, a ZeroMQ coordinator, EAGLE training
pipelines under FSDP2, asynchronous checkpointing via PyTorch DCP, and
SGLang as the rollout backend; the model-free fallback is Lookahead
{[}Paper §6.1; all code-confirmed, Phase 1{]}. Testbed: 8× DGX H100 = 64
GPUs (per node: 8 H100, 2× Intel Xeon Platinum 8480C, 2 TB RAM; 900 GB/s
NVLink, 400 Gb/s InfiniBand), plus a separate A100 cluster {[}Paper
§6.1{]}. Models: Qwen2.5-7B, DeepSeek-R1-Distill-Qwen-7B, Qwen2.5-32B,
Llama-3.3-70B-Instruct {[}Paper §6.1{]}. Algorithm: GRPO, temperature
0.9, max generation 32K, TP ∈ \{2,4,8\}, Adam, BF16 {[}Paper §6.1; Code
confirms n=8, KL \texttt{low\_var\_kl}/0.001, LR 1e-6{]}. Data:
Eurus-2-RL (math+code) for RL; an OpenThoughts2-1M subset to warm up the
drafter {[}Paper §6.1{]}. Baselines: Open-R1, VeRL (SOTA), TLT-Base
{[}Paper §6.1{]}. Metric: per-step token throughput averaged over three
steps after warm-up, i.e., gains are \emph{step-throughput at matched
reward curves} (Fig 12), not a claim of fewer steps to convergence
{[}Paper §6.1{]}.

\subsubsection{6.2 Consolidated results}\label{consolidated-results}

\begin{longtable}[]{@{}
  >{\raggedright\arraybackslash}p{(\linewidth - 6\tabcolsep) * \real{0.2500}}
  >{\raggedright\arraybackslash}p{(\linewidth - 6\tabcolsep) * \real{0.2500}}
  >{\raggedright\arraybackslash}p{(\linewidth - 6\tabcolsep) * \real{0.2500}}
  >{\raggedright\arraybackslash}p{(\linewidth - 6\tabcolsep) * \real{0.2500}}@{}}
\caption{Consolidated results: TLT vs.~baselines (per-step token
throughput; conditions and provenance in columns).
{[}§6.2{]}}\tabularnewline
\toprule\noalign{}
\begin{minipage}[b]{\linewidth}\raggedright
Result
\end{minipage} & \begin{minipage}[b]{\linewidth}\raggedright
Condition
\end{minipage} & \begin{minipage}[b]{\linewidth}\raggedright
Number
\end{minipage} & \begin{minipage}[b]{\linewidth}\raggedright
Source
\end{minipage} \\
\midrule\noalign{}
\endfirsthead
\toprule\noalign{}
\begin{minipage}[b]{\linewidth}\raggedright
Result
\end{minipage} & \begin{minipage}[b]{\linewidth}\raggedright
Condition
\end{minipage} & \begin{minipage}[b]{\linewidth}\raggedright
Number
\end{minipage} & \begin{minipage}[b]{\linewidth}\raggedright
Source
\end{minipage} \\
\midrule\noalign{}
\endhead
\bottomrule\noalign{}
\endlastfoot
End-to-end vs VeRL (headline) & over SOTA & \textgreater1.7× (abstract)
/ 1.7--2.1× (figure) & {[}Paper Abstract; Fig 11{]} \\
End-to-end geomean & H100 / A100 & 1.86× / 1.70× & {[}Paper Fig 11{]} \\
End-to-end peak & Qwen-32B, H100, 8 nodes & 2.12× & {[}Paper Fig 11; Tbl
3{]} \\
TLT vs VeRL, H100 & Qwen-7B / DS-7B / Qwen-32B / Llama-70B & 1.76 / 1.71
/ 2.12 / 1.86× & {[}Paper Fig 11{]} \\
TLT-Base vs VeRL, H100 & same models & 1.42 / 1.38 / 1.54 / 1.31× &
{[}Paper Fig 11{]} \\
Rollout share of step & , & ≈85\% & {[}Paper §1; Fig 1a{]} \\
SD speedup, peak grid & Qwen-32B, BS=1, depth 12/verify 64, temp 0 &
3.66× & {[}Paper Fig 13b{]} \\
SD accept length, max & depth 16/verify 128, temp 0 & 9.78 & {[}Paper
Fig 13a{]} \\
topK insensitivity & Qwen-32B, depth 12/verify 64 & 3.47--3.65× (topK
4--16) & {[}Paper Tbl 1{]} \\
Batch sensitivity (verify 64) & Qwen-32B, depth 10 & 3.62× (BS=1) →
1.70× (BS=32) & {[}Paper Tbl 4{]} \\
GPU diversity (rollout, BS1/TP1) & Qwen-7B micro-benchmark & B200 2.33×
· H100 2.61× · A100 2.79× · RTX5090 2.91× · RTX4090 2.87× · RTX3090
3.22× & {[}Paper Tbl 2{]} \\
Cluster scale (Qwen-7B) & 1/2/4/8 nodes & 1.21 / 1.45 / 1.62 / 1.76× &
{[}Paper Tbl 3{]} \\
Cluster scale (Qwen-32B) & 1/2/4/8 nodes & OOM / OOM / 1.83 / 2.12× &
{[}Paper Tbl 3{]} \\
Bucketed CUDAGraph & Llama-3-8B, TP=4, 4 strategies & 7.81 → 30.39 →
10.69 GB (2.8× vs naïve) & {[}Paper Tbl 5{]} \\
Selective async ckpt & drafter checkpoint & 893 → 280 → 97 ms (9.2×) &
{[}Paper Fig 17a{]} \\
Sequence packing & , & 13.3 → 29.6 samples/s (2.2×) & {[}Paper Fig
17b{]} \\
Adaptive drafter accept len & RL / downstream, Base→R & 4.59→6.53 /
3.76→5.15 & {[}Paper Tbl 6{]} \\
SD method (Eagle, default) & Qwen-7B, TP=2, BS=1 & 6.53 acc, 542.12
tok/s, 2.24×, 1× train cost & {[}Paper Tbl 7{]} \\
Case study (adaptive SD) & 128 req, Qwen-32B, H100 & 2.44× over no-SD &
{[}Paper Fig 14{]} \\
Overheads & , & stage-transition \textless1\%, SD reprefill \textless3
s, drafter every 10 steps & {[}Paper §6.4{]} \\
\end{longtable}

\subsubsection{6.3 GPU diversity and
scalability}\label{gpu-diversity-and-scalability}

SD's relative gain grows on more I/O-bound hardware: 2.33× on B200 up to
3.22× on RTX 3090 {[}Paper Tbl 2{]}. \emph{Caveat:} this is a BS=1, TP=1
rollout micro-benchmark, not full-scale training, it isolates the
memory→compute shift (slower memory subsystems gain more), and is not
representative of end-to-end training throughput. End-to-end gains grow
with cluster and model size: Qwen-7B 1.21×→1.76× (1→8 nodes); Qwen-32B
is OOM at 1--2 nodes and reaches 2.12× at 8 nodes {[}Paper Tbl 3{]},
larger jobs have longer tails and bigger bubbles to harvest. (See
\hyperref[figures]{Figures}: \texttt{tbl2\_gpu\_diversity},
\texttt{tbl3\_cluster\_scaling}, \texttt{fig11\_endtoend\_speedup}.)

\subsubsection{6.4 Drafter quality and SD-method
choice}\label{drafter-quality-and-sd-method-choice}

Among SD methods in TLT (Qwen-7B, TP=2, BS=1): Base/No-SD 242.11 tok/s
(1.00×); Eagle 542.12 tok/s, 2.24×, 1× cost; HASS 553.91 tok/s, 2.29×,
3× cost; Eagle-3 617.42 tok/s, 2.55×, 7× cost {[}Paper Tbl 7{]}. TLT
defaults to Eagle because higher-accept rivals require multi-pass
``training-time test'' that inflates per-step cost and convergence
difficulty, unaffordable in a small rollout bubble {[}Paper §6.5{]}.
OSD-style training further lifts a Qwen2.5-0.5B drafter (5.95→6.68→6.89)
and an Eagle drafter (1.57→6.53→6.77) {[}Paper Tbl 8{]}.

\subsubsection{6.5 Overheads and
losslessness}\label{overheads-and-losslessness}

Overheads are small: stage-transition \textless1\%; SD-switch reprefill
\textless3 s; drafter trained only every ≈10 steps {[}Paper §6.4{]}.
Most important for the thesis: VeRL and TLT reward curves overlap
closely for Qwen2.5-7B and -32B {[}Paper Fig 12{]}, direct evidence the
overlap-based speedup leaves learning dynamics unchanged. The gains are
\emph{utilization} (step throughput), not approximation; the paper does
not claim, and this survey does not assert, a reduction in the
\emph{number of steps} to a target reward.

\begin{center}\rule{0.5\linewidth}{0.5pt}\end{center}

\subsection{7. Why Not Fully Async?}\label{why-not-fully-async}

\subsubsection{7.1 The tempting next step, and its hidden
cost}\label{the-tempting-next-step-and-its-hidden-cost}

If idle GPUs can train the drafter, why not also update the
\emph{policy} asynchronously from partial rollouts during the long
decode? The paper addresses this in ``Can We Break RL Synchronization?''
{[}Paper §7{]}. Asynchronous RL with continuous batching ``mixes
multiple policy versions'' within a single rollout {[}Paper §7{]},
biasing gradient estimation and drifting on-policy → off-policy. TLT's
SD approach composes with \emph{limited} async RL ``without introducing
additional loss or policy staleness'' {[}Paper §7{]}.

\subsubsection{7.2 The lossless/off-policy contrast in the
model}\label{the-losslessoff-policy-contrast-in-the-model}

The §5.4 condition \(\eta_d I_t>\delta_\pi\) makes the tradeoff legible
{[}Author synthesis{]}. TLT tolerates \emph{drafter} staleness because
the drafter affects only SD's \emph{speed}, never the verified output
distribution, a stale drafter lowers the accept rate but cannot corrupt
a token (SD is token-level lossless, §2.2). Async RL instead tolerates
\emph{policy} staleness, which enters the gradient directly: there the
analogue of \(\delta_\pi\) is not a tracking error to be corrected but a
bias injected into the update. TLT confines all tolerated staleness to
the \emph{auxiliary} path and keeps the \emph{core} path exact.

\subsection{8. Implementation \&
Reproducibility}\label{implementation-reproducibility}

Grounded in the on-disk repository (\texttt{mit-han-lab/fastrl},
Apache-2.0, \texttt{main} @ \texttt{bce3df7}; full audit in
\texttt{phase1\_code\_grounding.md}).

Stack \& layout. Top-level \texttt{verl/} (281 .py, RL trainer, workers,
coordinator, DataBuffer), \texttt{eagle-train/} (12 .py, drafter
trainers + warm-up), \texttt{third-party/sglang/} (vendored in-tree,
1287 .py, adaptive SD, BEG-MAB, tree drafting, n-gram),
\texttt{examples/} (\texttt{bench\_sd.sh}, \texttt{grpo\_7B.sh},
\texttt{grpo\_32B\_multi\_nodes.sh}). Adaptive-SD work tracks SGLang PR
\#4732 (paper ref {[}59{]}). Released drafters:
\texttt{mit-han-lab/Qwen2.5-7B-Eagle-RL}, \texttt{-32B-Eagle-RL}
{[}Code: bench\_sd.sh:2; grpo\_*.sh{]}.

Reproduced run configs. 7B: TP=4, FSDP2, Ulysses SP=4, GRPO n=8, batch
64 prompts/mini 4, max-gen 32K, temp 0.9, KL
\texttt{low\_var\_kl}/0.001, \texttt{speculative.bs\_threshold=32}, 8
GPUs/1 node {[}Code: grpo\_7B.sh{]}. 32B: TP=8, SP=8, 8 nodes × 8 GPUs =
64, max-prompt 2048 {[}Code: grpo\_32B\_multi\_nodes.sh{]}, consistent
with the Tbl 3 8-node / 2.12× and OOM rows.

Key reproducibility caveat (the most important finding). The online
``drafter trained for free during RL'' pipeline is off by default in the
shipped config: \texttt{enable\_drafter\_training=false} and
\texttt{collect\_hidden\_states\_from\_sgl=false} {[}Code:
fastrl\_trainer.yaml:19,26{]}. The example scripts demonstrate adaptive
SD with a pre-trained \texttt{*-Eagle-RL} drafter; reproducing the
headline ``free byproduct'' claim requires enabling these flags. The
machinery (\texttt{EagleBackgroundTrainer} + coordinator + DataBuffer)
is fully present and wired.

Reception \& follow-on work (search-sourced, June 2026, post-cutoff).
TLT was covered by MIT News, \emph{``New method could increase LLM
training efficiency''} (2026-02-26), syndicated across MIT EECS /
Schwarzman College / Climate Portal and Mirage News; the article frames
the speedup in popular-press terms (\textasciitilde70--210\%), a loose
rendering of the paper's precise 1.7--2.1×, not adopted here.
Near-simultaneous concurrent work, \emph{``Beat the long tail:
Distribution-Aware Speculative Decoding for RL Training''}
(arXiv:2511.13841), addresses the same problem from a distribution-aware
angle; it is concurrent, not a TLT result.

\begin{center}\rule{0.5\linewidth}{0.5pt}\end{center}

\subsection{9. Threats to Validity / Assumptions
Ledger}\label{threats-to-validity-assumptions-ledger}

Dashboard assumptions (modeled, not paper-validated): 1. Pareto tail
(§3.2), paper Fig 1a is empirical/log-normal; the closed form is for
tractability. The dashboard's ``5.7× vs median'' is numerically wrong
(true ≈14×, §3.2). 2. Uniform per-token accept rate α (§5.1),
contradicted by Fig 16 (α decays with depth); the Markov model is
qualitative. 3. Exponential benefit decay, τ=16 (§5.3), a fit, not in
the paper; the deployed rule is a hard \texttt{≤32} step with 10-check
warmup. 4. Fisher information ≈ constant; local α≈e\^{}(−D\_KL) (§5.4),
convergence framing only. Plus the dashboard's §3 batch numbers (``1.08×
at BS=32'') are inconsistent with Tbl 4 and are not used.

Schematic / approximate figures: \texttt{fig1\_longtail\_*} and
\texttt{fig16\_*} are SCHEMATIC (only stated anchors are real);
\texttt{fig13} interior cells and \texttt{fig11} Open-R1 bars are
figure-read/approximate; \texttt{dashboard\_*} surfaces are modeled, not
measured (figures/manifest.md).

Code vs paper divergences (within-CONFIRMED): worker states are 5 not 3
(§4.2); ``frozen LM head'' is EAGLE-only, not EAGLE-3 (§4.4); DataBuffer
uses random sampling rather than explicit long-sequence prioritization
(§4.1); the 9.2× checkpoint figure is an end-to-end measurement, not a
parameter ratio (§4.3).

Scope of the speedup claims: all speedups are per-step token throughput
(or per-decode SD speedup), measured at matched reward curves {[}Paper
Fig 12{]}, not wall-clock-to-convergence or accuracy gains.
GPU-diversity numbers are a BS=1/TP=1 micro-benchmark. End-to-end gains
are workload-dependent: they grow with tail severity, model size, and
cluster size, and shrink when tails are short. The drafter is ``adequate
for SD speedup,'' not claimed SOTA as a standalone model; the BEG-MAB
carries ε=0.1 exploration overhead and assumes near-stationarity within
its 1000-step window.

\begin{center}\rule{0.5\linewidth}{0.5pt}\end{center}

\subsection{10. Related Work}\label{related-work}

The paper situates TLT against systems that relax synchronization
differently, and it is worth keeping the categories distinct (Phase 2
reconciliation):

\begin{itemize}
\tightlist
\item
  Off-policy / staleness-tolerant: AReaL (large-scale asynchronous RL)
  {[}arXiv:2505.24298{]} and APRIL (active partial rollouts that
  over-provision then recycle incomplete rollouts across steps)
  {[}arXiv:2509.18521{]} mix or carry over policy versions.
\item
  Disaggregation: StreamRL (disaggregated streaming generation)
  {[}arXiv:2504.15930{]}.
\item
  Synchronous long-tail scheduling: RollPacker (tail batching into
  long/short rounds) is itself synchronous {[}arXiv:2509.21009{]}, it is
  a scheduling method, \emph{not} an async-RL system; the paper notes
  its tail batching is limited because trajectories for identical
  queries vary widely in length {[}Paper §8{]}.
\end{itemize}

These accelerate training but (for the off-policy cases) risk degrading
quality by admitting stale responses or truncating rollouts. TLT's
differentiator is that it preserves the original RL algorithm and is
\emph{orthogonal}, it can layer on top of limited-async RL rather than
compete with it {[}Paper §7, §8{]}. The same overlap discipline
generalizes across RL algorithms: GRPO, RLOO, DAPO, REINFORCE,
REINFORCE++ share the workflow TLT accelerates {[}Paper §7{]}.

\begin{center}\rule{0.5\linewidth}{0.5pt}\end{center}

\subsection{11. Conclusion}\label{conclusion}

Future directions: extensions and open problems. The companion dashboard
proposes three forward-looking variants, design sketches \emph{not}
evaluated by the paper, inheriting the §5 assumptions.

The dashboard proposes three forward-looking variants, {[}Dashboard §5 /
Author synthesis{]}, design sketches \emph{not} evaluated by the paper,
inheriting the §5 assumptions.

A. Variance-aware thresholding. \(k^*(t)=k_0(1+\lambda\,V(t)/V(0))\),
\(V(t)=\mathrm{Var}(L_i\mid L_i>t)\) {[}Dashboard §5A{]}.
\emph{Rationale {[}Author synthesis{]}:} a high-variance residual tail
warrants enabling SD earlier, generalizing the static ``\textless32''
rule into a workload-adaptive one.

B. Multi-drafter ensembling.
\(\mathbb{E}[n_j]=\frac{1-\alpha_j^{\gamma+1}}{1-\alpha_j}\),
\(\mathbb{E}[n_{\text{ens}}]=\sum_{j} w_j\mathbb{E}[n_j]\),
\(w_j=|B_j|/B\) {[}Dashboard §5B{]}. \emph{Rationale:} different batch
regimes have different accept profiles; a weighted ensemble could
dominate any single drafter, at the memory cost Bucketed CUDAGraph
(§4.5) is built to contain.

C. Curriculum-ordered rollouts.
\(\sigma^*=\mathrm{argsort}_i(-\hat L_i/c_i(s))\), \(c_i\) = drafter
confidence {[}Dashboard §5C{]}. \emph{Rationale:} starting the
likely-longest, hardest-to-draft sequences first could shorten the
critical path.

Open problems {[}Dashboard §6{]}. (1) sample complexity of online
distillation vs the policy drift rate (the quantitative §5.4 condition);
(2) optimal depth/verify under \emph{non-uniform} accept
(\(\alpha_k=\alpha_1\beta^{k-1}\), relaxing §5.1); (3) a learnable
threshold trained by a secondary RL loop (replacing fixed
τ/``\textless32''); (4) cross-task drafter transfer bounds. The paper
independently notes compatible regimes, uniformly-long responses,
multi-turn tool-calling RL, fixed-target serving/edge deployment
{[}Paper §7{]}.

\begin{center}\rule{0.5\linewidth}{0.5pt}\end{center}

Concluding remarks. TLT's lesson is methodological as much as empirical:
the long-tail tax in reasoning RL is best paid back by overlapping
concurrent work in time, not by approximating the algorithm. By training
a one-layer drafter opportunistically on idle GPUs (non-blocking Spot
Trainer + DataBuffer), harvesting bubbles with a preemptible
coordinator, hiding checkpoint I/O (9.2×) and padding waste (2.2×), and
switching on token-level-lossless speculative decoding only in the
memory-bound tail (2.44× in the case study, peak 3.66× per-decode), TLT
delivers \textgreater1.7× end-to-end (1.7--2.1× over VeRL; geomean
1.86×/1.70× on H100/A100; peak 2.12×) {[}Paper Abstract, Fig 11, Tbl
3{]} while leaving reward curves indistinguishable from the synchronous
baseline {[}Paper Fig 12{]}. The mathematical model explains the
magnitude and saturation of these gains \emph{qualitatively}, with the
caveats that those closed forms are the dashboard author's
assumption-based scaffolding (one of which, the 5.7× Pareto example, is
numerically wrong) and that the released code refines several paper
details (5-state machine, EAGLE-only frozen LM head, online-training off
by default). The standing result: async-style utilization and
synchronous-grade correctness are not mutually exclusive, a design point
the fully-asynchronous RL literature has yet to occupy.

\begin{center}\rule{0.5\linewidth}{0.5pt}\end{center}

\subsection*{Acknowledgements}\label{acknowledgements}
\addcontentsline{toc}{subsection}{Acknowledgements}

This is an independent analysis. We gratefully acknowledge Hu, Yang,
Guo, Yao, Lin, Gu, Cai, Gan, Klimovic, and Han for the TLT/FastRL work
(ASPLOS '26, arXiv:2511.16665v2, CC-BY-4.0) and the released
\texttt{mit-han-lab/fastrl} code (Apache-2.0); all original empirical
results are theirs. Analytical models are adapted from the interactive
companion dashboard by Eric Hibbs.

\subsection*{Figures}\label{figures}
\addcontentsline{toc}{subsection}{Figures}

All figures are regenerated by \texttt{make\_figures.py} (vector
\texttt{.pdf}+\texttt{.svg}, colorblind-safe); provenance and
exact-vs-schematic status are in \texttt{figures/manifest.md}.

\begin{figure}
\centering
\includegraphics[width=0.92\linewidth,height=\textheight,keepaspectratio,alt={End-to-end RL training speedup vs VeRL=1.0, H100 + A100. {[}Paper Fig 11{]} (TLT/TLT-Base/VeRL exact; Open-R1 bars figure-read/approximate; some A100 Open-R1 cells illegible → n/a).}]{figures/fig11_endtoend_speedup.pdf}
\caption{End-to-end RL training speedup vs VeRL=1.0, H100 + A100.
{[}Paper Fig 11{]} (TLT/TLT-Base/VeRL exact; Open-R1 bars
figure-read/approximate; some A100 Open-R1 cells illegible → n/a).}
\end{figure}

\begin{figure}
\centering
\includegraphics[width=0.92\linewidth,height=\textheight,keepaspectratio,alt={SD speedup vs Tokens\_to\_Verify across batch sizes; the collapse to 1.70× motivates the elastic threshold. {[}Paper Tbl 4{]} (exact).}]{figures/tbl4_batch_erosion.pdf}
\caption{SD speedup vs Tokens\_to\_Verify across batch sizes; the
collapse to 1.70× motivates the elastic threshold. {[}Paper Tbl 4{]}
(exact).}
\end{figure}

\begin{figure}
\centering
\includegraphics[width=0.92\linewidth,height=\textheight,keepaspectratio,alt={Rollout SD speedup across GPUs (BS=1, TP=1 micro-benchmark). {[}Paper Tbl 2{]} (exact).}]{figures/tbl2_gpu_diversity.pdf}
\caption{Rollout SD speedup across GPUs (BS=1, TP=1 micro-benchmark).
{[}Paper Tbl 2{]} (exact).}
\end{figure}

\begin{figure}
\centering
\includegraphics[width=0.92\linewidth,height=\textheight,keepaspectratio,alt={End-to-end speedup vs cluster size; 32B OOM at 1--2 nodes. {[}Paper Tbl 3{]} (exact).}]{figures/tbl3_cluster_scaling.pdf}
\caption{End-to-end speedup vs cluster size; 32B OOM at 1--2 nodes.
{[}Paper Tbl 3{]} (exact).}
\end{figure}

\begin{figure}
\centering
\includegraphics[width=0.92\linewidth,height=\textheight,keepaspectratio,alt={Selective async checkpointing (893→280→97 ms, 9.2×) and sequence packing (13.3→29.6 samples/s, 2.2×). {[}Paper Fig 17{]} (exact).}]{figures/fig17_spot_trainer.pdf}
\caption{Selective async checkpointing (893→280→97 ms, 9.2×) and
sequence packing (13.3→29.6 samples/s, 2.2×). {[}Paper Fig 17{]}
(exact).}
\end{figure}

\begin{figure}
\centering
\includegraphics[width=0.92\linewidth,height=\textheight,keepaspectratio,alt={Accept length Target-Base→Target-R (+42\% RL). {[}Paper Tbl 6{]} (exact).}]{figures/tbl6_adaptive_drafter.pdf}
\caption{Accept length Target-Base→Target-R (+42\% RL). {[}Paper Tbl
6{]} (exact).}
\end{figure}

\begin{figure}
\centering
\includegraphics[width=0.92\linewidth,height=\textheight,keepaspectratio,alt={Depth×Verify grids for accept length and speedup; boxed extrema (9.78, 3.66×) verified. {[}Paper Fig 13{]} (boxed extrema exact; interior cells figure-read/approximate).}]{figures/fig13_sd_grid_heatmaps.pdf}
\caption{Depth×Verify grids for accept length and speedup; boxed extrema
(9.78, 3.66×) verified. {[}Paper Fig 13{]} (boxed extrema exact;
interior cells figure-read/approximate).}
\end{figure}

\begin{figure}
\centering
\includegraphics[width=0.92\linewidth,height=\textheight,keepaspectratio,alt={Accept rate vs draft-token index, adaptive vs vanilla. {[}Paper Fig 16, SCHEMATIC; endpoints 32\%/6\% from Dashboard{]}.}]{figures/fig16_accept_decay_SCHEMATIC.pdf}
\caption{Accept rate vs draft-token index, adaptive vs vanilla. {[}Paper
Fig 16, SCHEMATIC; endpoints \textsubscript{32\%/}6\% from
Dashboard{]}.}
\end{figure}

\begin{figure}
\centering
\includegraphics[width=0.92\linewidth,height=\textheight,keepaspectratio,alt={Modeled SD speedup over (α, γ). {[}Dashboard §2, assumption-based (uniform α); NOT measured{]}.}]{figures/dashboard_markov_surface.pdf}
\caption{Modeled SD speedup over (α, γ). {[}Dashboard §2,
assumption-based (uniform α); NOT measured{]}.}
\end{figure}

\begin{figure}
\centering
\includegraphics[width=0.92\linewidth,height=\textheight,keepaspectratio,alt={Aggregate SD benefit k·b(k), peak at k=τ=16. {[}Dashboard §3, assumption-based; NOT measured{]}.}]{figures/dashboard_threshold_curve.pdf}
\caption{Aggregate SD benefit k·b(k), peak at k=τ=16. {[}Dashboard §3,
assumption-based; NOT measured{]}.}
\end{figure}

\begin{figure}
\centering
\includegraphics[width=0.92\linewidth,height=\textheight,keepaspectratio,alt={Long-tail length distribution + under-utilized zone. {[}Paper Fig 1a/Fig 2, SCHEMATIC, not to scale; raw data not published as numbers{]}.}]{figures/fig1_longtail_SCHEMATIC.pdf}
\caption{Long-tail length distribution + under-utilized zone. {[}Paper
Fig 1a/Fig 2, SCHEMATIC, not to scale; raw data not published as
numbers{]}.}
\end{figure}

\begin{center}\rule{0.5\linewidth}{0.5pt}\end{center}

\subsection*{References}\label{references}
\addcontentsline{toc}{subsection}{References}

\emph{Primary and secondary sources, then external works actually cited.
External bibliographic details follow the TLT paper's reference list;
arXiv/venue strings verified in Phase 5
(\texttt{phase5\_citations.md}).}

{[}Primary, Paper{]} Q. Hu, S. Yang, J. Guo, X. Yao, Y. Lin, Y. Gu, H.
Cai, C. Gan, A. Klimovic, S. Han. \emph{Taming the Long-Tail: Efficient
Reasoning RL Training with Adaptive Drafter.} ASPLOS '26, Pittsburgh,
PA, Mar 22--26 2026. arXiv:2511.16665v2 (21 Jan 2026). DOI
10.1145/3779212.3790231. Code (Apache-2.0):
github.com/mit-han-lab/fastrl.

{[}Secondary, Dashboard{]} E. Hibbs. \emph{Taming the Long-Tail, an
interactive companion} (``Model Math'' §1--6).
https://erichibbsgh.github.io/Taming-the-Long-Tail-an-interactive./
(accessed 2026-06-05).

{[}Code artifact{]} \texttt{mit-han-lab/fastrl}, Apache-2.0,
\texttt{main} @ commit \texttt{bce3df7}; HF drafters
\texttt{mit-han-lab/Qwen2.5-\{7B,32B\}-Eagle-RL}.

External works cited: - Z. Shao et al.~\emph{DeepSeekMath} (GRPO).
arXiv:2402.03300, 2024. - DeepSeek-AI et al.~\emph{DeepSeek-R1.}
arXiv:2501.12948, 2025. - Y. Leviathan, M. Kalman, Y. Matias. \emph{Fast
Inference from Transformers via Speculative Decoding.} ICML 2023. - C.
Chen et al.~\emph{Accelerating LLM Decoding with Speculative Sampling.}
arXiv:2302.01318, 2023. - Y. Li, F. Wei, C. Zhang, H. Zhang.
\emph{EAGLE.} ICML 2024. · \emph{EAGLE-2.} EMNLP 2024. · \emph{EAGLE-3.}
arXiv:2503.01840, 2025. - L. Zhang, X. Wang, Y. Huang, R. Xu.
\emph{HASS: Learning Harmonized Representations for Speculative
Sampling.} ICLR 2025. - T. Cai et al.~\emph{Medusa.} arXiv:2401.10774,
2024. · Z. Ankner et al.~\emph{Hydra.} COLM 2024. - X. Miao et
al.~\emph{SpecInfer.} ASPLOS 2024. - X. Liu et al.~\emph{Online
Speculative Decoding (OSD).} ICML 2024. - G. Sheng et
al.~\emph{HybridFlow (VeRL).} EuroSys 2025. - Hugging Face.
\emph{Open-R1.} 2025. - W. Fu et al.~\emph{AReaL: A Large-Scale
Asynchronous RL System for Language Reasoning.} arXiv:2505.24298, 2025.
- Y. Zhong et al.~\emph{StreamRL.} arXiv:2504.15930, 2025. - Y. Zhou et
al.~\emph{APRIL: Active Partial Rollouts in RL to Tame Long-tail
Generation.} arXiv:2509.18521, 2025. - W. Gao et al.~\emph{RollPacker:
Mitigating Long-Tail Rollouts for Fast, Synchronous RL Post-Training.}
arXiv:2509.21009, 2025. - Y. Zhong et al.~\emph{RLHFuse (Optimizing RLHF
with Stage Fusion).} NSDI 2025. - Z. Mei et al.~\emph{ReaLHF.}
arXiv:2406.14088, 2024. - L. Zheng et al.~\emph{SGLang.} NeurIPS 2024. ·
P. Moritz et al.~\emph{Ray.} OSDI 2018. · P. Hintjens. \emph{ZeroMQ.}
O'Reilly, 2013. - PyTorch Team. \emph{FSDP2}; \emph{Torch Distributed
Checkpoint.} 2025. - M. M. Krell et al.~\emph{Efficient Sequence
Packing.} arXiv:2107.02027, 2022. - Y. Zhao et al.~\emph{Lookahead.} KDD
2024. - Qwen et al.~\emph{Qwen2.5 Technical Report.} arXiv:2412.15115,
2025. - G. Cui et al.~\emph{Eurus-2-RL / Process Reinforcement through
Implicit Rewards.} arXiv:2502.01456, 2025. - OpenThoughts.
\emph{OpenThoughts2-1M.} 2025. · L. Zheng et al.~\emph{MT-Bench.} 2023.
- Q. Yu et al.~\emph{DAPO.} arXiv:2503.14476, 2025. · A. Ahmadian et
al.~\emph{RLOO (Back to Basics).} ACL 2024. · J. Hu et
al.~\emph{REINFORCE++.} arXiv:2501.03262, 2025. - D. Kingma, J. Ba.
\emph{Adam.} ICLR 2015. · ByteDance. \emph{W\&B trace of DAPO
Reproduction on verl.} 2025.

Concurrent / reception (search-sourced, post-cutoff, dated): - MIT News.
\emph{New method could increase LLM training efficiency.}
news.mit.edu/2026/\ldots-0226 (2026-02-26). - \emph{Beat the long tail:
Distribution-Aware Speculative Decoding for RL Training.}
arXiv:2511.13841, 2025 (concurrent work).

\subsubsection{BibTeX for the primary
paper}\label{bibtex-for-the-primary-paper}

\begin{Shaded}
\begin{Highlighting}[]
\VariableTok{@inproceedings}\NormalTok{\{}\OtherTok{hu2026tlt}\NormalTok{,}
  \DataTypeTok{title}\NormalTok{     = \{Taming the Long{-}Tail: Efficient Reasoning RL Training with Adaptive Drafter\},}
  \DataTypeTok{author}\NormalTok{    = \{Hu, Qinghao and Yang, Shang and Guo, Junxian and Yao, Xiaozhe and}
\NormalTok{               Lin, Yujun and Gu, Yuxian and Cai, Han and Gan, Chuang and}
\NormalTok{               Klimovic, Ana and Han, Song\},}
  \DataTypeTok{booktitle}\NormalTok{ = \{Proceedings of the 31st ACM International Conference on Architectural}
\NormalTok{               Support for Programming Languages and Operating Systems, Volume 2 (ASPLOS \textquotesingle{}26)\},}
  \DataTypeTok{year}\NormalTok{      = \{2026\},}
  \DataTypeTok{doi}\NormalTok{       = \{10.1145/3779212.3790231\},}
  \DataTypeTok{note}\NormalTok{      = \{arXiv:2511.16665v2\}}
\NormalTok{\}}
\end{Highlighting}
\end{Shaded}

\emph{(The repo README's BibTeX reads ``31th'', corrected here to
``31st''; its News-line link text uses a non-canonical title, the title
above is the canonical PDF title.)}

\end{document}
